\chapter*{Abstract}
\addcontentsline{toc}{chapter}{Abstract}
Sequential quasi-Monte Carlo (SQMC) is an extension of sequential Monte Carlo (SMC) that propagates particles deterministically through low-discrepancy point sets, achieving an error rate smaller than the Monte Carlo rate of $O_P(N^{-1/2})$. However, SQMC is computationally intensive at $\mathcal{O}(N \log N)$ and suffers from the curse of dimensionality, since importance sampling does not generalise well to high dimensions. We develop a high-performance implementation of SQMC on parallel accelerators (GPU, TPU and multi-core CPU) using \texttt{jax} and substantiate its advantage on a toy filtering model. We then show how Rao-Blackwellisation overcomes high-dimension limitations of SQMC in the context of football matches by reducing the effective dimension of the latent state. This particle approach does not, however, outperform a factorial state-space model with an extended Kalman filter (EKF) approximation, indicating that the EKF is a suitable approximation for this problem.

This work is part of a contribution to the open-source package \texttt{cuthbert}, a library for efficient state-space inference algorithms in \texttt{jax}.

\textbf{Keywords}: Sequential Monte Carlo, Sequential Quasi-Monte Carlo, Low-discrepancy point sets, Hilbert sort, Rao-Blackwellisation, Parallel accelerators, Football

% We then show how Rao-Blackwellisation overcomes these limitations for football matches, reducing the effective dimension so that SQMC is applied to a smaller subset of the latent state space. 

% improving the performance of SQMC's two core primitives---low-discrepancy point generation and Hilbert-curve ordering---
% This comes from performing a Hilbert sort and importance sampling operation at every iteration.

% Sequential quasi-Monte Carlo (SQMC) is a quasi-Monte Carlo analogue of sequential Monte Carlo (SMC) that propagates particles deterministically through low-discrepancy point sets, achieving an error rate smaller than the Monte Carlo rate $O_P(N^{-1/2})$ and, unlike SMC, uniform-in-time error control. This thesis develops a high-performance implementation of SQMC on parallel accelerators (GPUs, TPUs, and multi-core CPUs) using JAX, exploiting the fact that SQMC's two core primitives---low-discrepancy point generation and Hilbert-curve ordering---are bit-level, embarrassingly parallel operations that map naturally onto accelerator hardware. We validate the implementation on a toy filtering model, demonstrating both the statistical advantages of SQMC over SMC and the throughput gains of accelerator execution. We then apply SQMC to football match modelling, where we perform Rao-Blackwellization to integrate out a subset of the latent variables analytically and run SQMC on the smaller remaining subset of parameters, substantially reducing the effective dimension of the filtering problem.